% Shared between the expanded supplementary source and official submission appendix.
% Update only this file when decoder/scoring/baseline protocol changes.
\section{Decoder and Training Protocol}
\label{sec:appendix-protocol}

\paragraph{Teacher and student modules.}
The teacher is a fixed TIGER T5 encoder--autoregressive-decoder model.
Students retain its frozen history encoder and SID embeddings.
A trainable context reader processes the encoder's history states:
SID-MLP uses attention, GRU uses recurrent readout, and CNN uses
temporal convolution and pooled features. All three feed four
prefix-conditioned \emph{MLP} SID prediction heads. The reported
trainable budgets (1.664M/2.304M/1.745M) include readers
\emph{and heads}, excluding the frozen encoder and embeddings.

\paragraph{Raw conditionals, valid paths, and complete scores.}
Digit CE and teacher-forced token KD operate on raw position-wise
256-token codebook distributions. Student candidate scores apply a
legal-extension mask and renormalize positions 2--4; the first
position uses the raw codebook softmax. Student serving and
candidate-scoring paths share this convention. HServe uses a
softmax over the finite teacher-proposed set ($M=64$ at training)
of cached teacher complete-item scores and corresponding student
valid-path scores. Its depth marginals sum \emph{these same}
candidate probabilities over shared-prefix groups.
DA-CDRS instead uses four per-position teacher log-probability
increments and four valid-path student increments to form local
gaps. The teacher cache's completed-candidate score may contain
an EOS contribution beyond four SID increments; equality of
the two teacher score conventions requires that remainder
to be constant across candidates of a request.
On-policy GKD/DistiLLM-2 teacher conditionals are normalized
within the selected 256-token codebook at each SID position.

\paragraph{Objectives and candidate sampling.}
The Full objective is digit CE plus $4$ times HServe and
$31.42$ times DA increment matching, without direct Pair Loss.
HServe uses four unit-weighted prefix KL terms from the shared
complete-candidate distribution, yielding effective conditional
KL weights $(4,3,2,1)$ by the chain rule. Temperature $\tau$
scales supported candidate scores before softmax and multiplies
the resulting HServe KL by $\tau^2$.
DA forms teacher-ranked cross-cutoff pairs between positions
1--10 and 11--40; by default all eligible pairs are used.
Its pairwise Softplus evaluation
$\operatorname{softplus}(0.05-G_S)$ determines \emph{detached}
weights, while valid-pair-averaged normalized Huber losses
supply the gradient. The independent Plain Pair ablation
directly optimizes the same Softplus instead of DA.
The sensitivity sweep's $\beta=31.42$ equals
$\lambda_D\lambda_{\mathrm{corr}}$ with $\lambda_D=1$.

\paragraph{Adapted KD comparisons.}
RD uses rank-decayed teacher preferences and dynamic student
ranks; RRD samples catalog negatives; RCE-KD matches candidate
rankings with cross-entropy. GKD mixes teacher-forced and
student-sampled SID-prefix supervision. DistiLLM-2 combines
skew-divergence matching with contrastive sampled SIDs:
teacher labels are detached, but sampled student scores
remain differentiable. SmartGR combines token-level and
beam-prefix objectives. The external KD entrypoint uses
a common $0.3$ digit CE plus $0.7$ token-KL base before each
objective-specific term; the Full method instead uses digit CE
as its base. Each method selects its checkpoint by
validation NDCG@10. These are objective adaptations on the
same student backbones, not reimplementations of every
original method's full architecture.
